\chapter{Transactions and State Transitions}

\section{Purpose}

\begin{principlebox}
For review purposes, a transaction is an authorization-bearing, metered request to apply one or more state transitions.
\end{principlebox}

That definition is deliberately broader than ``a signed payment.'' Some transactions transfer native assets. Some call contracts. Some deploy code. Some carry multiple instructions or messages. Some submit proofs. Some vote. Some update validator or committee state. Some are generated by system processes rather than ordinary users. Some fail but still consume fees, nonces, queue positions, or operational attention.

Chapter 1 gave us the state-machine model. Chapter 2 taught us to identify assets and authority. Chapter 3 studies the normal carrier of authority into the state machine: the transaction or transaction-like input.

The central mistake is to treat the transaction object as the outcome. A transaction is a request. The outcome is the effect produced when that request is interpreted by the current state, the execution rules, the resource limits, the ordering environment, and the settlement assumptions.

A serious transaction review asks:

\begin{codeblock}
Who authorized this request?
What state can it read or write?
What replay scope limits it?
What resources does it consume?
What can run before or after it?
What happens if it fails?
Which evidence proves the result?
When is the result safe to rely on?
\end{codeblock}

This chapter is not a full survey of Bitcoin, Ethereum, Solana, Sui, or Cosmos internals. Part III will do that. Here the goal is to learn a portable transaction-review method.

\section{Transactions Are Inputs, Not Outcomes}

\subsection{Transaction Object, Attempted Transition, Committed Effect}

Security review should separate three things: the transaction object, the attempted transition, and the committed effect.

The transaction object is the data submitted to the system. It may contain signatures, account references, calldata, script witnesses, message bodies, fee fields, gas limits, blockhashes, sequence numbers, chain identifiers, expirations, or object references.

The attempted transition is what the execution rules try to do with that object in the current state. The same object can behave differently if the state changes before inclusion. A swap transaction that was safe when constructed may become unsafe after pool reserves move. A withdrawal message may be safe after finality but unsafe before finality. A transaction that simulates successfully may fail if another transaction consumes the relevant state first.

The committed effect is what actually remains in canonical state after execution and settlement. It may include storage updates, balance changes, consumed nonces, emitted logs, fee payment, object version changes, queue updates, or no application state change at all if execution fails.

\begin{table}[htbp]
\centering
\small
\renewcommand{\arraystretch}{1.18}
\begin{tabularx}{\textwidth}{@{}>{\RaggedRight\arraybackslash}p{0.23\textwidth}>{\RaggedRight\arraybackslash}p{0.34\textwidth}X@{}}
\toprule
Layer & What to inspect & Common mistake \\
\midrule
Transaction object & Signed bytes, fields, targets, accounts, fees, replay scope & Assuming the human-facing description matches the signed payload \\
Attempted transition & Code path, state reads, state writes, resource use, failure behavior & Assuming a transaction will run against the state seen during construction \\
Committed effect & Canonical state, consumed authority, logs, receipts, finality status & Treating inclusion, event emission, or receipt success as proof of safety \\
\bottomrule
\end{tabularx}
\caption{Transaction review separates the request from the resulting state change.}
\end{table}

A valid transaction can still be unsafe. Validity means the system accepted the request under explicit rules. It does not mean the user understood the action, the price was fair, the oracle was fresh, the bridge source event was final, the protocol remained solvent, or the application encoded the intended invariant.

\subsection{Transaction Success Is Layered}

The phrase ``the transaction went through'' is too imprecise. It may mean that a wallet constructed it, a signer authorized it, an RPC accepted it, a mempool propagated it, a block included it, execution did not revert, a receipt was produced, an event was emitted, an indexer saw it, or settlement risk is low enough for the application.

Those are different facts. A bridge should not mint because a user says a transaction was broadcast. A protocol should not trust an event without checking the emitting contract and settlement assumption. A UI should not treat a simulation as a guarantee about actual inclusion.

Precision is a security habit. If the claim is about inclusion, say inclusion. If the claim is about successful top-level execution, say receipt success. If the claim is about durable cross-domain settlement, say which finality or settlement threshold is being used.

\section{Transaction Anatomy Across Models}

\subsection{The Same Review Questions Reappear}

Different systems package transactions differently. Bitcoin-style systems consume and create UTXOs. Ethereum-style systems use account transactions with nonces, calldata, gas, and signatures. Solana transactions contain account lists and instructions. Object systems such as Sui name objects and versions. Cosmos SDK applications route messages through an application pipeline.

The packaging differs, but the review questions repeat: who authorized the request, what state is touched, what replay boundary applies, what resources are consumed, what failure semantics apply, and what evidence confirms the result.

\begin{table}[htbp]
\centering
\small
\renewcommand{\arraystretch}{1.16}
\begin{tabularx}{\textwidth}{@{}>{\RaggedRight\arraybackslash}p{0.15\textwidth}>{\RaggedRight\arraybackslash}p{0.39\textwidth}X@{}}
\toprule
Model & What the transaction names & Security review focus \\
\midrule
UTXO & Inputs to consume and outputs to create & Which outputs are spent, which outputs are created, what script or witness authorizes the spend, and what value disappears as fee? \\
Account & Sender, nonce, target, value, calldata, gas fields, and signature & Which account authorizes, which target executes, which replay boundary applies, and what state or logs change? \\
Instruction & Signatures, account list, recent blockhash, and program instructions & Which accounts are writable, which program owns them, which signers are required, and which instruction fails? \\
Object & Input objects, versions, gas object, authenticator, and effects & Which objects are read or mutated, and which capability permits the change? \\
Message pipeline & Signed envelope containing one or more application messages & Which module handles the message, and which checks are admission-only versus deterministic execution? \\
\bottomrule
\end{tabularx}
\caption{Chapter 3 keeps execution-model detail shallow and focuses on portable review questions.}
\end{table}

Bitcoin's developer guide describes transactions as inputs and outputs, where each input spends a previous output and each output waits as a UTXO until later spent.\footnote{Bitcoin Developer Guide, \href{https://developer.bitcoin.org/devguide/transactions.html}{``Transactions''}.} Ethereum documentation describes transactions as signed instructions from accounts that update network state.\footnote{\href{https://ethereum.org/developers/docs/transactions/}{ethereum.org, ``Transactions''}.} Solana documentation describes transactions as containing instructions, account signatures, and a recent blockhash, with all instructions processed together.\footnote{Solana Foundation, \href{https://solana.com/docs/core/transactions}{``Transactions''}.} Cosmos SDK documentation separates transaction creation, mempool checks, delivery, and commit.\footnote{Cosmos SDK v0.47, \href{https://docs.cosmos.network/sdk/v0.47/learn/beginner/tx-lifecycle}{``Transaction Lifecycle''}.}

Those references matter, but this chapter should not make the reader memorize five transaction formats. The first-principles method is to locate the authority, state, replay scope, resources, failure behavior, evidence, and settlement assumption in whatever model is being reviewed.

\section{Lifecycle: From Construction to Settlement}

\subsection{A Transaction Changes Meaning Over Time}

A transaction lives through stages. Each stage answers a different security question.

\begin{figure}[htbp]
\centering
\begin{tikzpicture}[
  node distance=9mm and 8mm,
  compact node/.style={security node, text width=25mm, minimum height=7mm}
]
\node[compact node] (constructed) {Constructed};
\node[compact node, right=of constructed] (authorized) {Authorized};
\node[compact node, right=of authorized] (submitted) {Submitted};
\node[compact node, right=of submitted] (pending) {Pending};
\node[compact node, below=of pending] (included) {Included};
\node[compact node, left=of included] (executed) {Executed};
\node[compact node, left=of executed] (confirmed) {Confirmed};
\node[compact node, left=of confirmed] (settled) {Settled};

\draw[security arrow] (constructed) -- (authorized);
\draw[security arrow] (authorized) -- (submitted);
\draw[security arrow] (submitted) -- (pending);
\draw[security arrow] (pending) -- (included);
\draw[security arrow] (included) -- (executed);
\draw[security arrow] (executed) -- (confirmed);
\draw[security arrow] (confirmed) -- (settled);
\end{tikzpicture}
\caption{Transaction review should name the lifecycle stage instead of saying that the transaction ``went through.''}
\end{figure}

Construction is when the transaction is assembled. At this point, the target, parameters, accounts, amounts, deadlines, fee settings, and expected state may still be wrong. Wallets and frontends matter here because they shape what the user sees and signs.

Authorization is when an actor or rule permits the request. That may be a key signature, witness data, account validation, role check, module rule, capability, system process, or threshold signature. Authorization proves that the machine can accept the authority. It does not prove informed consent or safety.

Submission and pending propagation introduce ordering risk. A transaction may be visible to searchers, builders, validators, relayers, sequencers, or competing users. It can be copied, delayed, censored, replaced, sandwiched, or made stale.

Inclusion and execution answer whether a canonical block or batch contains the transaction and what execution did with it. Execution may succeed, revert, fail internally while the top-level transaction succeeds, consume fees, consume a nonce, emit logs, or update state.

Confirmation and settlement answer whether later systems should rely on the result. The required threshold differs for a wallet UI, exchange deposit, bridge mint, liquidation, rollup withdrawal, or governance action. Chapter 8 will treat finality and settlement in depth; here the lesson is simple: a transaction result is not equally safe for every downstream use at every lifecycle stage.

\section{Authorization, Replay, and Domain Scope}

\subsection{A Signature Is Not Intent}

A signature proves that a signing authority produced a valid signature over particular data under a particular scheme. It does not prove that the user understood the transaction, saw the real target, intended the same operation the protocol will execute, or was protected from phishing.

A user may believe they are signing a swap, while the signed payload grants an unlimited approval to an attacker. The signature can be valid and the transaction can execute successfully. The security failure is not cryptographic. The failure is that the user-side authority was tricked into authorizing the wrong transition.

Chapter 4 will go deeper into keys and signatures. Chapter 3 needs only the transaction-review rule:

\begin{codeblock}
Everything security-critical must be bound into the authorization:
chain or domain, signer, target, function, asset, amount, nonce,
deadline, and any intended constraints.
\end{codeblock}

If a field is not bound into what is authorized, an attacker may be able to move the authorization into a different context.

\subsection{Replay Is Reuse of Authority}

Replay is the reuse of valid authority where it should no longer be valid or should never have been valid. It can happen across time, chains, contracts, functions, accounts, forks, bridge domains, or deployments.

Replay protection is not one mechanism. Ethereum account nonces consume account-level transaction authority. Chain IDs bind transactions to a chain and prevent cross-chain replay in EIP-155-style signing.\footnote{Vitalik Buterin, \href{https://eips.ethereum.org/EIPS/eip-155}{``EIP-155: Simple Replay Attack Protection''}.} UTXO transactions consume specific outputs. Solana recent blockhashes bound transaction lifetime. Cosmos account sequences consume account-level authority. Bridges often need message nonces or nullifiers. Permit signatures need their own nonce and domain.

Do not say ``we have a nonce'' as if that solves replay. Ask what authority the nonce consumes. A transaction nonce may not protect a permit signature. A bridge message nonce may not protect a governance action. A withdrawal nullifier may not protect a cross-chain replay unless it is scoped to the right source, destination, message, and recipient.

The review prompt is:

\begin{codeblock}
What authority is consumed?
Who owns the nonce, sequence, object, output, or nullifier?
Is it scoped to the signer, chain, contract, function, asset, and amount?
Can it be reused after upgrade, fork, migration, or cross-chain delivery?
Does it expire?
\end{codeblock}

Replay protection must match the asset or authority being consumed.

\section{Fees, Ordering, and Failure Semantics}

\subsection{Resource Limits Are Security Boundaries}

Public blockchains cannot allow unbounded computation, storage growth, signature verification, account access, bandwidth use, or block space consumption. Fees, gas, compute budgets, size limits, account limits, and block limits are security boundaries, not merely UX annoyances.

On Ethereum, gas bounds execution. Ethereum's documentation describes gas as the fee required for transactions and contract execution.\footnote{\href{https://ethereum.org/developers/docs/gas/}{ethereum.org, ``Gas and Fees''}.} Solana uses base fees, optional prioritization fees, transaction size limits, and compute limits.\footnote{Solana Foundation, \href{https://solana.com/docs/core/fees}{``Fees''}.} UTXO systems commonly express fees as the difference between input value and output value. Cosmos-style systems distinguish gas, fee settings, mempool admission, ante handling, and deterministic delivery.

The model-specific details differ, but the security questions are stable. Can the required transition remain executable as state grows? Can an attacker make an operation too expensive? Who pays for failure? Can a queue, proof, liquidation, withdrawal, or governance action exceed a block or transaction limit? Are relayers, keepers, or users economically incentivized to submit the transaction when fees spike?

Resource failure is often a liveness failure before it becomes a direct theft. A protocol may still record that every user has a withdrawal claim, while no user can practically execute the withdrawal path.

\subsection{Ordering Is Part of the Transaction's Meaning}

Ordering changes outcomes. Consider two valid transactions:

\begin{codeblock}
tx_A: update oracle price
tx_B: borrow against collateral
\end{codeblock}

If \texttt{tx\_B} executes before \texttt{tx\_A}, the borrow sees the old price. If it executes after \texttt{tx\_A}, it sees the new price. Both orders may be valid under chain rules. Only one may be safe under the application's assumptions.

This is why pending transactions, mempools, private relays, builders, validators, sequencers, relayers, and keepers matter to application security. The execution path depends not only on the transaction fields but also on what can be observed, inserted before it, inserted after it, delayed, or censored.

The Part II chapters on peer-to-peer networks, block production, mempools, and MEV will go deeper. The Chapter 3 rule is enough for now: every transaction that depends on mutable state should state what must not change before inclusion and who may influence ordering.

\subsection{Failure Is an Outcome}

Failure is not absence of behavior. A failed transaction may consume fees. It may consume a nonce. It may leave a receipt. It may revert application state while preserving gas effects. It may block a queue if failure is handled poorly. It may make a keeper unprofitable. It may be exactly what an attacker wanted.

Atomicity is model-specific. In an EVM-style system, top-level revert rolls back state changes from that execution, but gas is still consumed and internal calls can fail if the caller catches or ignores the failure. In Solana, documentation describes instructions in a transaction as processed together, with failure of one instruction failing the entire transaction.\footnote{Solana Foundation, \href{https://solana.com/docs/core/transactions}{``Transactions''}.} Multi-transaction workflows, cross-chain workflows, and off-chain processes do not inherit single-transaction atomicity automatically.

For failure semantics, ask:

\begin{codeblock}
What exactly reverts?
What persists?
Who pays?
Can internal failure be swallowed?
Can one failing item block a batch?
Can repeated failure exhaust resources?
\end{codeblock}

Sometimes the safest transaction design is not the one that tries to do everything atomically. It is the one that makes partial progress explicit, bounded, observable, and recoverable.

\section{Evidence: State, Receipts, Logs, Events, and Simulation}

\subsection{Evidence Is Not Authority}

Observers rarely see state transitions directly. They use evidence: receipts, logs, events, traces, RPC responses, indexers, simulations, and UI interpretations. Each is useful. None should be mistaken for the security property itself.

\begin{table}[htbp]
\centering
\small
\renewcommand{\arraystretch}{1.18}
\begin{tabularx}{\textwidth}{@{}>{\RaggedRight\arraybackslash}p{0.20\textwidth}>{\RaggedRight\arraybackslash}p{0.28\textwidth}X@{}}
\toprule
Evidence & Useful for & Warning \\
\midrule
State & Authority for future validity and balances & Must be read from the correct canonical state \\
Receipt & Execution status, gas used, logs, block context & Does not prove economic safety or user intent \\
Event or log & Off-chain discovery, monitoring, analytics, bridge watchers & Communicates; does not define authority by itself \\
Indexer view & Search, UI, accounting, automation & Can lag, omit reorgs, or interpret incorrectly \\
Simulation & Preflight, UX warnings, keeper checks & Runs against assumed state and ordering, not future reality \\
\bottomrule
\end{tabularx}
\caption{Transaction evidence must be interpreted according to what produced it and what it can prove.}
\end{table}

Ethereum receipts include execution information such as block context, gas used, logs, and a status field for post-Byzantium transactions.\footnote{\href{https://ethereum.org/developers/docs/apis/json-rpc/}{ethereum.org, ``JSON-RPC API''}.} A receipt success means top-level execution did not revert. It does not mean the oracle price was fair, the user understood the signature, the trade was not sandwiched, the bridge source state was final, or the protocol remained solvent.

Events and logs are communication. They are useful for indexers, UIs, monitoring, analytics, and cross-domain watchers. But an event can be emitted without the expected state change if the code is wrong. A state change can happen without the expected event. An event can be read from a block that later disappears from canonical history. A watcher must verify the emitting contract, the event semantics, the proof or RPC source, and the finality assumption.

Simulation answers a narrower question: what would happen if this transaction ran now against this assumed state and context? It does not answer what will happen when the transaction is actually included, ordered, executed, and settled.

The practical rule is short:

\begin{codeblock}
Events are evidence.
Receipts are evidence.
Simulations are evidence.
Canonical state under the right settlement assumption is authority.
\end{codeblock}

\section{Worked Example: Swap Transaction Review}

Consider a user signing a swap through a decentralized exchange frontend:

\begin{codeblock}
swapExactTokensForTokens(
    tokenIn = USDC,
    tokenOut = WETH,
    amountIn = 10,000 USDC,
    minAmountOut = 3.8 WETH,
    path = [USDC, WETH],
    recipient = user,
    deadline = now + 10 minutes
)
\end{codeblock}

The user-facing summary says: swap 10,000 USDC for at least 3.8 WETH. That summary is useful, but the review has to inspect the transaction path.

\begin{table}[htbp]
\centering
\small
\renewcommand{\arraystretch}{1.18}
\begin{tabularx}{\textwidth}{@{}>{\RaggedRight\arraybackslash}p{0.20\textwidth}>{\RaggedRight\arraybackslash}p{0.29\textwidth}X@{}}
\toprule
Review area & Evidence & Security question \\
\midrule
Authorization & User signature, allowance, router target & Did the user authorize this router, path, recipient, amount, deadline, and chain? \\
Replay scope & Nonce, chain ID, deadline, transaction target & Can the authorization be reused in another time, contract, chain, or context? \\
State touched & Pool reserves, token balances, allowance state & Which balances, reserves, approvals, and accounting entries change? \\
Ordering & Pending visibility, public mempool or private route & Who can see, delay, sandwich, or reorder the swap? \\
Resources & Gas limit, fee settings, token behavior & Can execution fail or become stuck because resources are insufficient? \\
Failure & Revert on low output, token transfer behavior & If the swap fails, what fees are paid and does the approval remain? \\
Evidence & Receipt, events, balance checks, settlement & Do logs match state, and is the block safe enough for the next action? \\
\bottomrule
\end{tabularx}
\caption{A swap review traces authority, state, ordering, resources, failure, evidence, and settlement.}
\end{table}

The most important safety control is not merely that the user signed. The signature proves authorization over some payload. The safety controls are the minimum output, deadline, target contract, token assumptions, replay scope, slippage bound, and the user's ability to understand what is being authorized.

Now compare a malicious transaction such as \texttt{approve(attacker, unlimited)}. It may be well encoded, signed by the user, included in a valid block, executed successfully, and accompanied by an approval event. It is still unsafe for the user. The failure is that user authority was converted into a harmful transition.

This is the chapter's core lesson. Transaction review is not satisfied by asking whether the transaction is valid. It asks whether the resulting state transition is safe under the intended security model.

\section*{Review Questions}

\begin{enumerate}
\item What is the difference between a transaction object, an attempted transition, and a committed effect?
\item Why is ``the transaction was signed'' not enough to prove user intent or protocol safety?
\item What does replay protection consume, and why must it match the authority being protected?
\item Why can transaction inclusion, receipt success, event emission, and settlement mean different things?
\item Give an example of a transaction that succeeds but leaves an unsafe state.
\item Why are gas, fees, compute units, size limits, and block limits security concerns?
\item Why is simulation useful but insufficient as evidence of future execution?
\item Why should bridge watchers and indexers treat events as evidence rather than authority?
\end{enumerate}

\section*{Applied Exercises}

\begin{enumerate}
\item \textbf{Transaction anatomy.} Compare a Bitcoin UTXO spend, an Ethereum ERC-20 transfer, and a Solana transaction with two instructions. For each, identify the authority, replay boundary, state consumed or written, resource constraint, failure behavior, and evidence of success.

\item \textbf{Replay scope.} A protocol accepts an off-chain signature for a withdrawal, but the signed message includes only \texttt{user} and \texttt{amount}. Identify possible replay attacks, then propose the fields and state updates needed to make the authorization single-use and domain-scoped.

\item \textbf{Event versus state.} A bridge watcher sees a \texttt{Deposit(user, amount, nonce)} event and mints wrapped tokens on another chain. List what the watcher must verify besides the event data, including emitting contract, source-chain settlement, replay protection, and destination-chain state.

\item \textbf{Valid but unsafe.} For each case, explain why the transaction may be valid and still unsafe: an unlimited token approval to a malicious spender; a borrow after one-block oracle manipulation; a governance execution that transfers treasury assets after satisfying quorum; a withdrawal batch that fails because one recipient reverts; a swap that succeeds after being sandwiched.
\end{enumerate}
